\par\noindent\textbf{(1)} 势能 \(U=\rho g\int_{-a}^{a} y\sqrt{1+y'^2}\,dx\)，长度约束为 \(\int\sqrt{1+y'^2}\,dx=\ell\)。乘子法得到泛函 \(F=(\rho gy+\Lambda)\sqrt{1+y'^2}\)。由于不显含 \(x\)，Beltrami积分 \(F-y'F_{y'}=(\rho gy+\Lambda)/\sqrt{1+y'^2}=H\) 为常量，\(H\) 即水平张力。
\par\noindent\textbf{(2)} 设 \(c=H/(\rho g)\)，整理 \(1+y'^2=[(y+\Lambda/\rho g)/c]^2\)，积分得对称解 \(y(x)=c\cosh(x/c)+y_0\)。弧长元 \(ds=\cosh(x/c)\,dx\)，所以 \(\ell=2c\sinh(a/c)\)。在 \(\ell>2a\) 时这一定义 \(c>0\) 的唯一对称解。
\par\noindent\textbf{(3)} 索上张力沿切向，水平分量固定为 \(H=\rho gc\)；因切线斜率 \(y'=\sinh(x/c)\)，总张力 \(T(x)=H\sqrt{1+y'^2}=\rho gc\cosh(x/c)\)。最低点处最小，支点张力为 \(\rho gc\cosh(a/c)\)，而竖向总支持力 \(2H\sinh(a/c)=\rho g\ell\)。
\par\noindent\textbf{(4)} 在小横向位移 \(\eta(s,t)\) 且忽略纵向耦合的局部短波近似，线性式为 \(\rho\eta_{tt}=\partial_s[T(s)\eta_s]\)，所以局部波速 \(v(s)=\sqrt{T(s)/\rho}=\sqrt{gc\cosh(x/c)}\)。当 \(\ell\downarrow2a\) 时 \(c\to\infty\)，张力发散，说明完全绷直的极限不能用有限张力静态近似描述。
